\documentclass[12pt]{article}
\usepackage{graphicx} %package to manage images
\graphicspath{ {./figures/} }
\usepackage{caption}
\usepackage[font=scriptsize]{subcaption}
\captionsetup[figure]{labelsep=none}
\captionsetup[table]{labelsep=none}
\usepackage{bbm}
\usepackage{amsmath}
\usepackage{import}
\usepackage{array}
\usepackage{booktabs}     % no extra vertical space between paragraphs
\usepackage{indentfirst}
\usepackage{afterpage}
\usepackage{floatrow}
\usepackage{pdflscape}
\usepackage{soul}
\usepackage{float}
\usepackage{adjustbox}
\usepackage{longtable}
\usepackage{caption}
\usepackage{setspace}
\usepackage{afterpage}
\usepackage[margin=1in]{geometry}
\usepackage[round]{natbib}
\usepackage{hyperref}
\usepackage{titlesec}
\usepackage{threeparttable}
\usepackage{setspace}
\usepackage{float}
\floatstyle{plaintop}
\restylefloat{table}
\usepackage{rotating}
\usepackage{ragged2e}
\usepackage{authblk}

% Optional styling for nicer author/affil typography
\renewcommand\Authfont{\normalsize\bfseries}
\renewcommand\Affilfont{\small\normalfont}
\setlength{\affilsep}{0.3em} % space between author and affiliation

% \author[1,2,3]{Diane Coffey\thanks{Corresponding author: coffey@utexas.edu}}

% \author[1]{Siddhant Pandit}

% \affil[1]{Population Research Center, UT Austin}
% \affil[2]{Department of Sociology, UT Austin}
% \affil[3]{r.i.c.e., a research institute for compassionate economics}

\setlength{\parindent}{0pt}   % standard indent (~0.25 inch)
\setlength{\parskip}{1em}       % no extra vertical space between paragraphs

\titleformat{\section}
  {\normalfont\bfseries\fontsize{14}{16}\selectfont}  % Bold 14pt
  {}{0pt}{}  % No numbering

\titleformat{\subsection}
  {\normalfont\bfseries\fontsize{12}{14}\selectfont}  % Bold 12pt
  {}{0pt}{}  % No numbering

\titleformat{\subsubsection}
  {\normalfont\itshape\fontsize{12}{14}\selectfont}  % Italic 12pt
  {}{0pt}{}  % No numbering
  
\begin{document}
\RaggedRight
 
\title{Fertility, Household Wealth, and Maternal Nutrition among Indian Social Groups} 
\date{\today} 
\maketitle
   
% \begin{abstract}
% Maternal undernutrition is a threat to women’s health and a leading cause of child mortality and morbidity in India. This paper uses nationally representative data to estimate the fraction of women who begin pregnancy underweight -- a key indicator of maternal nutrition -- and presents the first estimates by caste and religion.  We use nonparametric reweighting to show that more than one in five Indian women start pregnancy underweight.  Adivasis and Dalits are 55\%, and 40\%, more likely, respectively, to be underweight before pregnancy than Forward Caste Hindus. Differences in fertility levels and birth spacing do not explain the maternal nutrition gaps that we observe.  Instead, differences in household wealth explain a large fraction of the gaps. Notably, and despite multidimensional disadvantage, we find that a similar fraction of Muslim and Forward Caste Hindu women start pregnancy underweight.  Our results suggest the continued need for monitoring of and health interventions in the prepregnancy period.
% \end{abstract}

% \subsection*{\textbf{Keywords:} Maternal nutrition, caste, pregnancy, decomposition}

\subsection*{Highlights}
{\setstretch{0.25}
\begin{itemize}
    \item The fraction of women starting pregnancy with BMI$<18.5$ was 42\% in 2005 and 22\% in 2020.
    \item Pregnant Adivasi and Dalit women are most likely to be underweight.
    \item Fertility and birth spacing \emph{do not} explain gaps across social groups.
    \item Differences in household wealth \emph{do} explain gaps across social groups.
\end{itemize}
}
\newpage
\doublespacing

\end{document}